\documentclass[aps,prx,reprint,superscriptaddress]{revtex4-2}

\usepackage{graphicx}
\usepackage{amsmath}
\usepackage{siunitx}
\usepackage{booktabs}

\begin{document}

\title{Supplemental Material: \\
   Quantum-referenced refractive-index calibration of air at 1762 nm for environmental phase compensation in quantum optical interfaces}

\author{Wei Wu}
\affiliation{Albert-Ludwigs-Universit\"at Freiburg, Physikalisches Institut, Hermann-Herder-Stra\ss e 3, 79104 Freiburg, Germany}
\affiliation{EUCOR Centre for Quantum Science and Quantum Computing, University of Freiburg, Hermann-Herder-Stra\ss e 3, 79104 Freiburg, Germany}
\author{Tobias Sch{\"a}tz}
\affiliation{Albert-Ludwigs-Universit\"at Freiburg, Physikalisches Institut, Hermann-Herder-Stra\ss e 3, 79104 Freiburg, Germany}
\affiliation{EUCOR Centre for Quantum Science and Quantum Computing, University of Freiburg, Hermann-Herder-Stra\ss e 3, 79104 Freiburg, Germany}
\author{Ulrich Warring}
\affiliation{Albert-Ludwigs-Universit\"at Freiburg, Physikalisches Institut, Hermann-Herder-Stra\ss e 3, 79104 Freiburg, Germany}

\date{\today}

\maketitle

\section{Campaign temporal coverage}
\label{sec:S0}

The processed dataset spans 231 calendar days (2025-01-01 to 2025-08-20) but is not a continuous record: it consists of 104 observation segments, separated by gaps longer than 2~h, with 89 active observation days in total. The longest gap (111 days) separates a sparse January--February block from the main body of the campaign in June--August, which contains 91.7\% of the observations and entirely exceeds the 25\,$^{\circ}$C validity bound of the Mathar model. The overall duty cycle, defined as the ratio of active segment time (35.6 days) to the calendar span, is 15.3\%. Within active periods the median sampling interval is 12.9~s.

Table~\ref{tab:S0} summarises the monthly coverage. January and February are the only months within the Mathar validity domain ($T \in [10,25]\,^{\circ}$C); February is represented by only 25 retained observations spread over 13 active days. The retained-observation mean conditions ($\overline{T} = 28.9\,^{\circ}$C, $\overline{\mathrm{RH}} = 30.2\%$, $\overline{P} = 986.4$\,hPa) are weighted toward the warm-season segments. A time-weighted mean is not reported: the 111-day gap separates the January--February block from the June--August body of the campaign, so the calendar-duration weight of the two observations bordering the gap dominates any such average, making it unrepresentative of either the observed conditions or a continuous campaign. The regression coefficients of Table~\ref{tab:coefficients} are estimated from the retained observations.

\begin{table*}[htbp]
\centering
\caption{Monthly coverage of the processed dataset. $n_{\rm obs}$: retained observations; active days: calendar days containing observations; cadence: median intra-month sampling interval; gap: longest intra-month gap; ranges are min--max within each month. The last column is the fraction of the month's
observations within the Mathar validity domain ($T \in [10,25]\,^{\circ}$C).}
\label{tab:S0}
\begin{tabular}{lcccccccc}
\hline
Month & $n_{\rm obs}$ & Active & Cadence & Gap & $T$ range & RH range & $P$ range & In-domain \\
  & & days   & (s) & (d)  & ($^{\circ}$C) & (\%) & (hPa) & \\
\hline
2025-01 & 12\,109 & 17 & 13.4 & 4.5 & 17.5--21.4 & 28.5--40.9 & 964.6--1006.8 & 100\% \\
2025-02 & 25 & 13 & 41\,317 & 5.9 & 20.2--22.1 & 28.5--38.4 & 982.2--999.4 & 100\% \\
2025-06 & 28\,911 & 13 & 12.8 & 1.5 & 29.3--33.3 & 22.8--34.6 & 980.6--995.7 & 0\% \\
2025-07 & 54\,397 & 26 & 12.8 & 4.0 & 25.9--34.6 & 18.9--44.0 & 974.1--993.3 & 0\% \\
2025-08 & 50\,342 & 20 & 12.9 & 1.0 & 25.4--31.9 & 24.8--45.3 & 978.9--994.5 & 0\% \\
\hline
\multicolumn{2}{l}{Total} & 89 & 12.9 & 111.3 & 17.5--34.6 & 18.9--45.3 & 964.6--1006.8 & 8.3\% \\
\hline
\end{tabular}
\end{table*}

\section{Statistical validation of the environmental coefficients}
\label{sec:S1}

The refractive index time series exhibits serial correlation characteristic of slowly varying environmental parameters. The residual autocorrelation function decays to $1/e$ at a lag of approximately 78 samples, corresponding to a decorrelation time well resolved by the 10--20~s measurement cadence. To ensure that reported coefficient uncertainties account for this structure, we compute heteroskedasticity- and autocorrelation-consistent (HAC) standard errors with a Bartlett kernel \cite{newey1987hypothesis,andrews1991heteroskedasticity} using a maximum lag of 156 samples, twice the measured $1/e$ decorrelation time. As the primary uncertainty estimate we use moving-block bootstrap resampling with a block length of 156 samples and $B = 5000$ replications \cite{kunsch1989jackknife,politis1994stationary}; the block length preserves the serial correlation structure of the original data while rendering successive blocks approximately independent. The HAC standard errors serve as an independent cross-check: the HAC and bootstrap standard errors agree to within 3\% for all three coefficients (Table~\ref{tab:S1b}), confirming that the uncertainty estimates are robust to the specific autocorrelation-consistent estimator employed. For reference, a HAC estimate with an insufficient lag of 3 samples, which does not capture the dominant serial correlation, underestimates the standard errors by approximately a factor of five.

First-order residual autocorrelation is additionally addressed using the Prais-Winsten generalized least squares estimator \cite{prais1954trend}. The transformation improves the Durbin-Watson statistic from 1.197 (raw OLS residuals) to 2.226 (transformed residuals), where a value of 2.0 indicates no first-order autocorrelation \cite{durbin1950testing}. The Prais-Winsten coefficient estimates agree with the OLS estimates within \SI{1.5}{\percent}, confirming that the point estimates are insensitive to the residual correlation structure.

The pressure coefficient reported in Table~I of the main text ($+2.5949 \times 10^{-7}\,\text{hPa}^{-1}$) follows from the regression coefficient ($+2.595 \times 10^{-9}$ in the native fit unit of $\text{Pa}^{-1}$) by conversion of the pressure unit from Pa to hPa (factor $10^2$). The result agrees with first-principles Lorentz--Lorenz scaling of the air refractivity at standard conditions.

\subsection{Uncertainty budget for single-epoch refractive index}
\label{sec:S1a}

Table~\ref{tab:S1a} provides the uncertainty budget for the determination of the refractive index at a single measurement epoch. The dominant contribution is the residual scatter of the regression ($\delta n = 1.84 \times 10^{-7}$), which includes the spatial and temporal mismatch between the environmental sensor and the optical path as well as interferometric readout noise. A decomposition of this residual against a linear surrogate of the Mathar model evaluated on the identical data rows shows that the intrinsic nonlinearity of the environmental response contributes only $4.4 \times 10^{-8}$, with the remaining $1.78 \times 10^{-7}$ attributable to measurement noise. The \SI{780}{\nano\meter} frequency reference contributes $\delta n = 1.0 \times 10^{-9}$; this follows from locking the \SI{780}{\nano\meter} laser to the $^{87}$Rb D$_2$ line via Doppler-free saturated absorption spectroscopy ($\delta\nu \approx \SI{384}{\kilo\hertz}$ at $\nu \approx \SI{384.230}{\tera\hertz}$), corroborated by long-term wavemeter monitoring of the laser frequency. The \SI{1762}{\nano\meter} quantum frequency reference and interferometer non-linearity are negligible. The combined expanded uncertainty ($k = 2$) is $3.68 \times 10^{-7}$.

\begin{table}[htbp]
\centering
\caption{Uncertainty budget for single-epoch refractive index determination at \SI{1762}{\nano\meter}.}
\label{tab:S1a}
\begin{tabular}{lc}
\hline
Source & $\delta n$ (1$\sigma$) \\
\hline
Residual scatter (incl.\ spatial/temporal mismatch) & $1.84 \times 10^{-7}$ \\
\quad of which: measurement noise & $1.78 \times 10^{-7}$ \\
\quad of which: model nonlinearity & $4.4 \times 10^{-8}$ \\
\SI{780}{\nano\meter} frequency reference   & $1.0 \times 10^{-9}$ \\
\SI{1762}{\nano\meter} frequency reference & $<10^{-12}$ \\
Interferometer non-linearity & Negligible \\
\hline
Combined standard uncertainty & $1.84 \times 10^{-7}$ \\
Combined expanded uncertainty ($k = 2$)  & $3.68 \times 10^{-7}$ \\
\hline
\end{tabular}
\end{table}

\subsection{Uncertainty budget for environmental coefficients}
\label{sec:S1b}

Table~\ref{tab:S1b} separates the statistical and systematic contributions to the uncertainty of the environmental coefficients. The statistical component is the moving-block bootstrap standard error ($B = 5000$, block length 156 samples), cross-checked against the HAC estimate. The systematic component accounts for the sensor gain and offset errors, that is, the uncertainty in the multiplicative scale and in the additive zero of the sensor readings. In a single-channel linear fit a constant offset is absorbed by the intercept $n_0$ and does not affect the coefficient, whereas the gain error propagates directly into the coefficient estimate as $\alpha_X^{\rm meas} = \alpha_X^{\rm true}/g_X$, where $g_X$ is the sensor gain. In the full measurement chain the offset is retained as well: it shifts the conditions at which the Ciddor anchor is evaluated, so the joint envelope of Sec.~\ref{sec:S2} is taken over the gain and the offset together.

Worst-case gain bounds follow from the manufacturer specifications of the BME280 sensor over the observed campaign ranges ($T$: \SIrange{17.5}{34.6}{\degreeCelsius}, span 17.1 K; $H$: \SIrange{18.9}{45.3}{\percent}, span 26.4\%; $P$: \SIrange{964.6}{1006.8}{\hecto\pascal}, span 42.2 hPa). For temperature, the accuracy of \SI{1.0}{\degreeCelsius} bounds the gain error to $2 \times 1.0/17.1 = 11.7\%$. For humidity, the accuracy of \SI{3}{\percent} at \SI{25}{\degreeCelsius} degrades to about \SI{4}{\percent} over the temperature range of the campaign (together with long-term drift), bounding the gain error to $2 \times 4.0/26.4 = 30.3\%$. For pressure, the accuracy of \SI{1}{\hecto\pascal} bounds the gain error to $2 \times 1.0/42.2 = 4.7\%$. These worst-case gain bounds translate directly into the endpoint sensitivities of Table~\ref{tab:S1b}: the change in the coefficient if the gain deviated from unity by the full bound, $\Delta\alpha_X^{\rm end} = (\delta g/g)_{\max} \cdot |\alpha_X|$.

For the standard systematic uncertainty $u_{\rm syst}$ we convert the worst-case gain bound into a standard deviation assuming a rectangular distribution of the gain over $[1-\delta, 1+\delta]$, i.e. division by $\sqrt{3}$. The values listed in Table~\ref{tab:S1b} are obtained from the joint gain--offset envelope of Sec.~\ref{sec:S2}: the same manufacturer bounds are propagated through all three channels and maximised over the envelope that the specifications justify, and the resulting maximum displacement is divided by $\sqrt{3}$. Relative to the single-branch treatment, in which the worst-case bound is applied to that coefficient alone and converted by $\sqrt{3}$, this reduces the systematic standard uncertainty by factors of $70$, $3.9$ and $53$ for $\alpha_T$, $\alpha_H$ and $\alpha_P$, respectively. For the humidity channel the single-branch rectangular value alone would be $3.99 \times 10^{-9}/\sqrt{3} = 2.30 \times 10^{-9}$; it is retained here only as a diagnostic, because it rescales the humidity coefficient in isolation and omits the Ciddor-anchor and Mathar branches. The paired bootstrap standard error of the difference quoted in Sec.~\ref{sec:S2} ($u_{\rm stat} = 1.30 \times 10^{-9}$) is larger than the single-coefficient bootstrap error listed here because it accounts for the correlation between the two estimates.

\begin{table*}[htbp]
\centering
\caption{Uncertainty budget for the environmental coefficients at \SI{1762}{\nano\meter}. Statistical uncertainties are moving-block bootstrap standard errors ($B = 5000$, block length 156 samples); HAC (lag 156) standard errors agree within \SI{3}{\percent}. The systematic column $u_{\rm syst}$ is obtained from the joint gain--offset envelope of Sec.~\ref{sec:S2}: the BME280 manufacturer bounds on the gain and offset are propagated through all three channels and maximised over the joint envelope justified by the specifications, and the maximum displacement is converted to a standard uncertainty assuming a rectangular gain distribution (division by $\sqrt{3}$). The endpoint column is the single-branch coefficient change for a gain deviation of the full worst-case bound, listed for reference only; $u_{\rm tot}$ is the root-sum-square of $u_{\rm stat}$ and $u_{\rm syst}$.}
\label{tab:S1b}
\begin{tabular}{lccccccc}
\hline
Coefficient & Value & $u_{\rm stat}$ & Endpoint & $u_{\rm syst}$ & $u_{\rm tot}$ & $U$ ($k=2$) & HAC SE \\
\hline
$\alpha_T$ (\si{K^{-1}})   & $-8.8474\times 10^{-7}$ & $1.41\times 10^{-9}$ & $1.03\times 10^{-7}$ & $8.57\times 10^{-10}$ & $1.65\times 10^{-9}$ & $3.30\times 10^{-9}$ & $1.41\times 10^{-9}$ \\
$\alpha_H$ (\%RH$^{-1}$)   & $-1.3152\times 10^{-8}$ & $1.15\times 10^{-9}$ & $3.99\times 10^{-9}$ & $5.91\times 10^{-10}$ & $1.29\times 10^{-9}$ & $2.59\times 10^{-9}$ & $1.16\times 10^{-9}$ \\
$\alpha_P$ (\si{hPa^{-1}}) & $+2.5949\times 10^{-7}$ & $9.85\times 10^{-10}$ & $1.23\times 10^{-8}$ & $1.34\times 10^{-10}$ & $9.94\times 10^{-10}$ & $1.99\times 10^{-9}$ & $9.69\times 10^{-10}$ \\
\hline
\end{tabular}
\end{table*}

\section{Same-condition comparison with the Mathar model}
\label{sec:S2}

To compare the measured environmental response with the Mathar model on equal footing, we evaluate the Mathar formulation \cite{mathar2007refractive} at the identical $(T, H, P)$ rows of the measurement campaign and fit the same linear surrogate model $n = n_0 + \alpha_T T + \alpha_H H + \alpha_P P$ to both the measured refractive index and the Mathar predictions. Both fits therefore share the same design matrix, and the coefficient difference is estimated with a paired moving-block bootstrap (identical block resampling for both models, $B = 5000$, block length 156 samples) that accounts for the correlation between the two estimates.

Table~\ref{tab:S2} summarizes the result for two analysis windows: the full campaign and the subset within the validity domain of the Mathar model ($T \in [10, 25]\,\si{\degreeCelsius}$, comprising 12,134 of the 145,784 observations; 91.7\% of the data lie above \SI{25}{\degreeCelsius}). For the full campaign, the measured $\alpha_T$ and $\alpha_P$ agree with the Mathar surrogate within \SI{0.35}{\percent} and \SI{0.93}{\percent}, respectively, and the humidity coefficient is smaller in magnitude than the surrogate by $+24.8\%$. Within the Mathar validity domain the humidity difference is not significant, and this subset has limited discriminating power. Because a consistent propagation of the sensor gain and offset bounds over the three channels limits the effect on the humidity coefficient to $1.02 \times 10^{-9}$ (\SI{24}{\percent} of the difference), and because the significance of the difference depends on the block length used in the paired resampling, we do not attribute the humidity-coefficient difference to a wavelength-specific physical effect; it is reported as an unresolved systematic.

The significance of the humidity difference can be assessed in two ways that answer different questions. The endpoint sensitivity of Table~\ref{tab:S1b} ($3.99 \times 10^{-9}$) rescales only the fitted humidity coefficient and omits the Ciddor-anchor and Mathar branches; it is therefore a single-branch estimate rather than a bound on the coefficient difference. Propagating the same manufacturer bounds through all three channels and maximising over the joint gain--offset envelope that the specifications justify gives a maximum displacement of $1.02 \times 10^{-9}$ for the humidity coefficient, i.e. \SI{24}{\percent} of the full-campaign difference $\Delta\alpha_H = +4.33 \times 10^{-9}$. With the rectangular treatment this corresponds to $u_{\rm syst} = 5.9 \times 10^{-10}$; combined with the paired bootstrap standard error of the difference, $u_{\rm stat} = 1.30 \times 10^{-9}$, the observed difference is $3.0\, u_{\rm comb}$ with $u_{\rm comb} = \sqrt{1.30^2 + 0.59^2} \times 10^{-9} = 1.43 \times 10^{-9}$. The statistical significance is conditional on the resampling block length: $u_{\rm stat} = 1.13 \times 10^{-9}$ at $0.56$~h blocks, $3.69 \times 10^{-9}$ at $7.17$~h blocks and $5.38 \times 10^{-9}$ at $24$~h blocks, so zero is excluded only for the shortest block length. The gain term is therefore no longer the limiting systematic. The remaining difference is compatible with the unresolved slow drift of the fringe-ratio anchor: the data-driven monthly injection is $4.4 \times 10^{-9}$ and the shape-bounded magnitude is $13.4 \times 10^{-9}$, either of which can account for the whole difference. This is an uncertainty-budget statement only: reducing the gain sensitivity does not restore a physical claim. The humidity difference remains compatible with the unresolved $1/R(t)$ drift and with the model-surrogate limits.

\begin{table*}[htbp]
\centering
\caption{Same-condition comparison of the measured environmental coefficients with a linear surrogate of the Mathar model fitted on identical data rows.$\Delta\alpha = \alpha^{\rm data} - \alpha^{\rm Mathar-surrogate}$;intervals are 95\% bootstrap percentile intervals.}
\label{tab:S2}
\begin{tabular}{lcccc}
\hline
Analysis & $N$ & $\Delta\alpha_T$ & $\Delta\alpha_H$ & $\Delta\alpha_P$ \\
\hline
In-domain ($10$--$25\,\si{\degreeCelsius}$) & 12,134 & $-1.41\times 10^{-8}$ & $+4.62\times 10^{-9}$ & $+3.79\times 10^{-9}$ \\
\quad 95\% interval & & $[-5.51, +2.61]\times 10^{-8}$ & $[-6.86, +21.5]\times 10^{-9}$ & $[+0.10, +8.05]\times 10^{-9}$ \\
Full campaign  & 145,784 & $-3.11\times 10^{-9}$ & $+4.33\times 10^{-9}$ & $+2.39\times 10^{-9}$ \\
\quad 95\% interval & & $[-6.26, -0.15]\times 10^{-9}$ & $[+1.79, +6.88]\times 10^{-9}$ & $[+0.31, +4.45]\times 10^{-9}$ \\
\hline
\end{tabular}
\end{table*}

A blocked out-of-sample validation complements the coefficient comparison: the data are divided into contiguous time blocks of 156 samples and five folds; within each fold the linear surrogate is fitted on the training blocks (including the intercept) and evaluated on the held-out blocks. Identical folds are used for both models and no global preprocessing is applied before splitting. The out-of-sample root-mean-square residual of the data surrogate is $(1.84 \pm 0.02) \times 10^{-7}$, whereas that of the Mathar surrogate is $(4.4 \pm 0.2) \times 10^{-8}$. The larger residual of the data surrogate reflects measurement noise that is absent in the analytic Mathar function; the small surrogate residual confirms that the linear model captures the Mathar environmental response to high accuracy over the observed parameter range.

\section{Phase budget for a representative free-space path}
\label{sec:S3}

For a representative \SI{1}{\meter} free-space path at \SI{1762}{\nano\meter}, the accumulated phase error under diurnal environmental swings of amplitude $\Delta T = \SI{1}{\degreeCelsius}$, $\Delta H = \SI{12}{\percent}$, and $\Delta P = \SI{5}{\hecto\pascal}$ is computed as

\begin{equation}
\Delta\phi = \sqrt{(\alpha_\phi^T\Delta T)^2 + (\alpha_\phi^H\Delta H)^2 + (\alpha_\phi^P\Delta P)^2},
\end{equation}

where $\alpha_\phi^X \equiv \partial\phi/\partial X$ per meter (Table~II of the main text). The per-meter phase budget is summarized 
in Table~\ref{tab:S3}.

\begin{table*}[htbp]
\centering
\caption{Phase error budget for a \SI{1}{\meter} free-space optical path at \SI{1762}{\nano\meter} under typical diurnal environmental swings. Sensitivities are from Table~II of the main text.}
\label{tab:S3}
\begin{tabular}{lccc}
\hline
Parameter  & $\Delta X$ (diurnal)  & $\alpha_\phi^X$  & $\Delta\phi$ (\si{\radian}) \\
\hline
Temperature & \SI{1}{\degreeCelsius}  & \SI{3.16}{\radian\,m^{-1}\,K^{-1}}  & \num{3.16}  \\
Relative humidity  & \SI{12}{\percent}  & \SI{0.047}{\radian\,m^{-1}\,(\%RH)^{-1}}  & \num{0.56}  \\
Pressure   & \SI{5}{\hecto\pascal}  & \SI{0.93}{\radian\,m^{-1}\,hPa^{-1}}  & \num{4.63}  \\
\hline
Total (RSS) &  &  & \textbf{\num{5.63}} \\
\hline
\end{tabular}
\end{table*}

Feed-forward phase compensation uses real-time environmental sensor data to compute and cancel the predicted phase shift before each gate operation. The residual phase error after compensation has two sources: the predictive accuracy of the linear model, characterized by the regression residual standard deviation $\sigma_n = 1.84 \times 10^{-7}$, and the uncertainty in the coefficient estimates themselves, which propagates through the environmental swings $\Delta X_j$ via the coefficient standard errors $\sigma_{\alpha_j}$. The combined residual is

\begin{equation}
\delta\phi_{\text{res}} = \frac{2\pi L}{\lambda}\,
  \sqrt{\sigma_n^{2} 
  + (\sigma_{\alpha_T}\Delta T)^2 
  + (\sigma_{\alpha_H}\Delta H)^2 
  + (\sigma_{\alpha_P}\Delta P)^2}.
\end{equation}

Using the statistical (bootstrap) coefficient standard errors of Table~\ref{tab:S1b} and the diurnal swings of Table~\ref{tab:S3}, the coefficient-uncertainty term contributes $1.5 \times 10^{-8}$ in refractive-index units, an order of magnitude below $\sigma_n$. The residual is therefore dominated by the predictive accuracy of the model, and the combined expression evaluates to $\delta\phi_{\text{res}} = \SI{0.66}{\radian}$ per meter.

For a fair model comparison, the static offset between the Mathar prediction and the measured refractive index must be removed, since any deployed system removes this offset with a single calibration. The static offset over the campaign is $\delta n = 1.166 \times 10^{-6}$, corresponding to \SI{4.16}{\radian} per meter (\SI{17.46}{\radian} over the full \SI{4.2}{\meter} round-trip path). After subtracting this offset, the Mathar residual standard deviation is $1.93 \times 10^{-7}$ ($\SI{0.69}{\radian}$ per meter), compared with $1.84 \times 10^{-7}$ ($\SI{0.66}{\radian}$ per meter) for the empirical model. The improvement in residual amplitude is therefore $I_{\rm amp} = 4.86\%$ (95\% interval $[3.9, 6.0]\%$), and the improvement in residual variance is $I_{\rm var} = 9.5\%$ (95\% interval $[7.7, 11.5]\%$). These modest improvements reflect that the empirical model and the debiased Mathar model agree closely over the observed parameter range, consistent with the same-condition coefficient comparison of Sec.~\ref{sec:S2}.

\section{Model limitations and scope of the measurement approach}
\label{sec:S4}

Several factors set the achievable accuracy of the empirical model. First, the model assumes linear, decoupled environmental dependences; while cross-correlations among $T$, $H$, and $P$ are small in our well-ventilated measurement corridor, a full multivariate polynomial including cross-terms could further reduce $\sigma_n$ toward the $10^{-8}$ level, as suggested by the small nonlinearity component of the residual (Table~\ref{tab:S1a}). Second, CO$_2$ variations are not tracked in the present setup; seasonal CO$_2$ changes of \SIrange{15}{25}{ppm} alter $n$ by approximately $3 \times 10^{-8}$, which lies below the current residual standard deviation but may become relevant as sensor performance improves.

Beyond these model-level uncertainties, the measurement approach itself has inherent scope limitations that define what the calibrated coefficients can and cannot correct in a deployed system. (i)~\textit{Local gradients}: the environmental sensor samples temperature, humidity, and pressure at a single location, whereas the optical path extends over several meters. Under strong thermal stratification or localized humidity sources, the sensor reading may deviate from the path-averaged value, introducing an uncompensated residual that scales with the gradient magnitude. This is particularly relevant for vertical beam paths or links traversing multiple thermal zones. (ii)~\textit{Turbulence}: the sensor sampling rate (\SI{2}{\hertz}) is far slower than the characteristic timescales of turbulent refractive-index fluctuations (milliseconds to seconds) \cite{tatarskii1961wave}. The feed-forward scheme therefore corrects only the slowly varying background but not the stochastic Kolmogorov-phase noise, which must be addressed by other means such as high-bandwidth phase tracking or spatial mode filtering. (iii)~\textit{Unsensed geometric path-length variations}: the dual-wavelength common-mode rejection cancels mechanical fluctuations only to the extent that both wavelengths experience identical geometric displacement. Differential effects, such as chromatic beam wander due to thermal lensing in transmissive optics or wavelength-dependent thermal expansion of mechanical mounts, produce an uncompensated residual that is not sensed by the environmental monitor. (iv)~\textit{Fast mechanical noise}: the interferometer rejects common-mode vibrations within its fringe-tracking bandwidth, but spectral components above this bandwidth, or differential vibrations between the two wavelength channels due to spatially separated beam paths on the detection side, are not canceled and set a practical noise floor for the phase correction. These scope limitations are inherent to any single-point, slow-sampling feed-forward approach and apply equally to the empirical and the Mathar-based compensation.

\section{Evidence grading of the measurement-chain budget}
\label{sec:S5}

Tables~\ref{tab:S4a}--\ref{tab:S4c} grade every term of the measurement-chain budget as \emph{quantified}, \emph{bounded} or \emph{unresolved}, together with its assumed magnitude, its temporal behaviour, the evidence it rests on and its propagated effect. The notebook numbers in the evidence column refer to the analysis notebooks archived with the data release. The grading follows the chain-propagation protocol of Sec.~\ref{sec:S2} and is summarised in the evidence-grading figure of the main text. It supports two design conclusions: a tighter humidity calibration is no longer the limiting systematic (the joint gain--offset envelope contributes at most \SI{24}{\percent} of the humidity difference), and a better frequency reference would not materially improve the coefficient discrimination. The entries whose totals depend on the pending block-length decision and on the held-out and statistical repairs of the analysis notebooks are marked provisional.

% Evidence-grading tables for the NS10657 supplemental material.
% Generated from code/analysis/evidence_grading.json (producer: 11_evidence_grading.ipynb).
% Three floats, one per output; labels tab:S4a/b/c. Included by supplemental_material.tex.

\begin{table*}[htbp]
\centering
\caption{Evidence grading, Output~A: terms that determine the absolute refractive index $n_{1762}$. Each source is graded \emph{quantified}, \emph{bounded} or \emph{unresolved}, matching panel~(b) of the evidence-grading figure of the main text. ``Magnitude'' is the assumed or measured size of the term; ``propagated effect'' states how it enters the refractive-index determination. Terms absorbed by the fitted intercept do not enter the environmental coefficients.}
\label{tab:S4a}
\footnotesize
\setlength{\tabcolsep}{4pt}
\begin{tabular}{@{}p{0.15\textwidth}p{0.095\textwidth}p{0.185\textwidth}p{0.095\textwidth}p{0.14\textwidth}p{0.20\textwidth}@{}}
\hline
\textbf{Source} & \textbf{Grade} & \textbf{Magnitude} & \textbf{Temporal behaviour} & \textbf{Evidence} & \textbf{Propagated effect} \\
\hline
Anchoring / interferometric offset \emph{(hard boundary; not correctable with this dataset)} & unresolved & $1.166\times10^{-6}$ (= 17.46 rad at 4.2 m; 4.16 rad/m) & static over the campaign; no independent monitor & nb 02; SM S1a; hard-boundary note & limits the absolute scale and baseline and blocks any ``Mathar absolutely correct'' claim; absorbed by the fitted intercept, so it does not enter the coefficients \\
K convention ($f_{780}/f_{1762}$ vs documented 1762.17/780.24) \emph{(convention / rounding, not a measurement)} & bounded & $dK/K = -5.17$ ppm $\to$ offset $5.2\times10^{-6}$, 4.4$\times$ the current $1.166\times10^{-6}$; 0.01 nm rounding alone allows $\pm$6.4 ppm & static multiplicative & nb 10 sec.~D; $f_{\rm rb}/f_{1762} = 2.2584857037$ & intercept only: $dK/K$ from $10^{-9}$ to $10^{-7}$ leaves $d(\alpha_H)_{\rm read}$ unchanged to four digits \\
780 nm frequency reference \emph{(Rb D2 saturated-absorption lock)} & quantified & $dn = 1.0\times10^{-9}$ ($d\nu \sim 384$ kHz at 384.230 THz) & static / long-term, wavemeter-monitored & SM S1a & enters single-epoch $n$ only; negligible for the coefficients \\
1762 nm quantum frequency reference & quantified & $dn < 10^{-12}$ & static & SM S1a & negligible \\
Interferometer non-linearity \emph{(declared negligible)} & bounded & not quantified & static & SM S1a & negligible \\
Residual scatter (incl.\ spatial/temporal mismatch) \emph{(statistical; the decomposition is not source identification)} & quantified & $1.84\times10^{-7}$; of which measurement noise $1.78\times10^{-7}$ and model nonlinearity $4.4\times10^{-8}$ & short-term scatter plus a slow component & SM S1a; nb 03 & sets single-epoch precision; the spatial/temporal mismatch inside it remains unresolved \\
Sensor-to-path spatial mismatch \emph{(possible contribution)} & unresolved & not separable; folded into the $1.84\times10^{-7}$ residual & short-term and slow & SM S1a & may contribute to the coefficient difference; must not be quoted as resolved \\
Wavelength convention (Mathar 1.762 vs 1.762174 $\mu$m; anchor 780.24 vs 780.246 nm) \emph{(documentation item)} & quantified & $d(dH) = -0.21\times10^{-12}$ and $+0.57\times10^{-12}$; anchor $dn = -0.037$ ppb & static & nb 10 sec.~D & $\leq 10^{-12}$ on the coefficients: no recomputation needed, only explicit documentation \\
\hline
\end{tabular}
\end{table*}

\begin{table*}[htbp]
\centering
\caption{Evidence grading, Output~B: terms that move the humidity coefficient difference $\Delta\alpha_H$, matching panel~(a) of the evidence-grading figure of the main text. ``Magnitude'' is the assumed or measured size of the term; ``propagated effect'' states how it enters the coefficient difference.}
\label{tab:S4b}
\footnotesize
\setlength{\tabcolsep}{4pt}
\begin{tabular}{@{}p{0.15\textwidth}p{0.095\textwidth}p{0.185\textwidth}p{0.095\textwidth}p{0.14\textwidth}p{0.20\textwidth}@{}}
\hline
\textbf{Source} & \textbf{Grade} & \textbf{Magnitude} & \textbf{Temporal behaviour} & \textbf{Evidence} & \textbf{Propagated effect} \\
\hline
Statistical uncertainty of the humidity difference \emph{(conditional on block length; not robust across the resampling procedures examined)} & quantified & $u_{\rm stat} = 1.30\times10^{-9}$ (SM nominal); $1.13\times10^{-9}$ at 0.56 h blocks; $3.69\times10^{-9}$ at 7.17 h blocks & block-length dependent; the paired-bootstrap SE grows from $1.13\times10^{-9}$ (0.56 h) to $5.38\times10^{-9}$ (24 h) & SM S1b; nb 06/07; \texttt{paired\_diff\_block\_sensitivity.json} & zero is excluded only at 0.56 h blocks; 7.17 h and 24 h blocks include zero \\
Humidity sensor gain (H channel) \emph{(consistent three-branch propagation, joint $(g,b)$ envelope)} & bounded & SM single-branch endpoint $3.99\times10^{-9}$ (corrected coordinate) vs joint envelope max $1.024\times10^{-9}$ ($u_{\rm syst} = 5.91\times10^{-10}$) & assumed constant multiplicative; BME280 $\pm$3 \%RH at $25\,^\circ$C, $\sim$4 \%RH over the campaign plus drift & nb 09; \texttt{joint\_gain\_offset\_envelope}; BME280 specification & moves $d(\alpha_H)_{\rm read}$ by $\leq 1.02\times10^{-9} = 24\%$ of the difference; gain is no longer the limiting systematic (the SM quote overestimates by 3.9$\times$) \\
Temperature and pressure sensor gain \emph{(same joint envelope)} & bounded & T: $1.03\times10^{-7} \to 1.484\times10^{-9}$; P: $1.23\times10^{-8} \to 2.324\times10^{-10}$ & assumed constant & nb 09 & cross-channel coupling is real and the response is not additive: the joint maximum exceeds the triangle sum of the channel-only maxima (H 1.024 vs $0.952\times10^{-9}$; linearised bound $0.915\times10^{-9}$), so a joint envelope is required \\
Sensor offset $b$ (centred parameterisation) \emph{(feasible envelope only; the two tutorial scenarios leave it at the endpoint gain)} & bounded & $|b| \leq \delta$; envelope vertex $b = -0.645$ (H) / $-0.125$ (T) / $-0.033$ (P); uncentred $b = (g-1)X_0 = 9.09$ \%RH at $g = 1.303$, 2.27$\times$ the spec bound & static & nb 09; task-card sec.~1 & the uncentred and the centred-$b=0$ hypotheses are different assumptions and both exceed the specification at the endpoint gain; only envelope vertices are legal sensitivity slices \\
Slow $1/R(t)$ component of the fringe-ratio anchor \emph{(shape not identifiable; bounded magnitude)} & unresolved & observed slow component 0.147--0.411 ppm; leverage $2.6$--$32.5\times10^{-9}$ per ppm; injection $+4.37\times10^{-9}$ (data-driven), up to $13.4\times10^{-9}$ (shape bound) & slow (ramp / seasonal); corr$(T,t) = 0.743$; the 111-day gap separates the two blocks & nb 10; \texttt{one\_over\_R\_term.json} & can account for the entire $4.33\times10^{-9}$ difference; it cannot be corrected with this dataset \\
Mathar validity-domain coverage \emph{(model-surrogate limitation, independent of calibration quality)} & unresolved & 91.7\% of rows above $25\,^\circ$C; in-domain $N = 12{,}134$ ($\sim$1.8 day equivalent), January plus 25 February rows; 28 intra-run gaps $> 2$ h (max $\sim$145 h) & static coverage property & nb 05; \texttt{campaign\_time\_coverage.json}; SM S0 & the in-domain comparison has limited discriminating power; the full-campaign comparison is an extrapolation and must be labelled as such \\
Block length for the coefficient uncertainties \emph{(decision pending, Task 4)} & unresolved & nominal $L = 156$ samples; 200 samples proposed & sets the statistical uncertainty, not a physical source & nb 07; todo v4 Task 4 & determines whether the difference excludes zero; no single nominal SE should be quoted before the decision \\
K as a reconstruction rather than a calibration \emph{(consistency check, not an independent validation)} & quantified & relative spread $3.95\times10^{-13}$; ptp/mean $2.98\times10^{-12}$ & no structure at this spread & nb 08; \texttt{chain\_propagation.calibrate\_K} & confirms that the toolset reproduces the archived pipeline; it does not validate the convention \\
\hline
\end{tabular}
\end{table*}

\begin{table*}[htbp]
\centering
\caption{Evidence grading, Output~C: refractivity change and slow phase compensation. ``Magnitude'' is the assumed or measured size of the term; ``propagated effect'' states how it enters the applied result.}
\label{tab:S4c}
\footnotesize
\setlength{\tabcolsep}{4pt}
\begin{tabular}{@{}p{0.15\textwidth}p{0.095\textwidth}p{0.185\textwidth}p{0.095\textwidth}p{0.14\textwidth}p{0.20\textwidth}@{}}
\hline
\textbf{Source} & \textbf{Grade} & \textbf{Magnitude} & \textbf{Temporal behaviour} & \textbf{Evidence} & \textbf{Propagated effect} \\
\hline
Offset-fair predictive improvement \emph{(over the observed interval and bandwidth)} & quantified & $I_{\rm amplitude} = 4.86\%$ [3.92, 5.95]; $I_{\rm variance} = 9.47\%$ [7.68, 11.54]; $\sigma_{\rm emp} = 1.8367\times10^{-7}$ vs $\sigma_{\rm Mathar} = 1.9304\times10^{-7}$ & campaign-scale & nb 02 & the surviving applied result; it must stay conditional on the observation interval and bandwidth \\
Frequency reference (780 nm / 1762 nm / K convention) \emph{(negative result for the design question Q2)} & quantified & 780 nm $dn = 1.0\times10^{-9}$; 1762 nm QFR $dn < 10^{-12}$; the K convention is intercept-only & static & SM S1a; nb 08/10 & a better frequency reference would not materially improve coefficient discrimination or slow compensation \\
Slow $1/R(t)$ drift \emph{(as Output~B row 5)} & unresolved & see Output~B row 5 & slow & nb 10 & limits predictive performance over intervals longer than the decorrelation time; it needs independent drift monitoring \\
Environmental coverage and observation schedule \emph{(campaign property, not a sensor limit)} & bounded & duty cycle 15.3\%; 111-day gap; in-domain $\sim$1.8 day equivalent & campaign-scale & nb 05; \texttt{campaign\_time\_coverage.json} & performance statements must be conditional on the schedule and the environmental range \\
Sensor gain and offset \emph{(as Output~B rows 2--4)} & bounded & joint envelope max $1.024\times10^{-9}$ (H) & static & nb 09 & small relative to the drift term for compensation; not the binding constraint \\
\hline
\end{tabular}
\end{table*}


\bibliography{reference}
\end{document}